\documentstyle{report}
%\def{\parameter{#1}}{{\it #1}}

\begin{document}

\chapter{IDAnalyticBH}

\section{Introduction}

\subsection{Purpose of Thorn}

Thorn IDAnalyticBH provides analytic initial data for vacuum black 
hole spacetimes. Initial data is provided for the 3-metric, extrinsic
curvature, and if appropriate the conformal factor and it's spatial 
derivatives. The current initial data sets are for a single (Schwarzschild)
black hole in isotropic coordinates, up to four Brill-Lindquist black 
holes, and any number of Misner-type black holes.

\subsection{Technical Specification}

\begin{itemize}

\item{Implements} einsteinID
\item{Inherits from} einstein
\item{Tested with thorns} Einstein

\end{itemize}

\section{Theoretical Background}


\section{Algorithmic and Implementation Details}

This thorn uses no special numerical methods, however two points
are worth noting

\begin{enumerate}

\item{} The solution for Misner is obtained by summing a sequence

\item{} The spatial derivatives of the conformal metric (when required) 
        are calculated accurately using finite differencing of the 
        exact solution by a very small spacing

\end{enumerate}

\section{Using the Thorn}

This thorn can provide either the physical metric (use\_conformal=''no'')
or the conformal metric and a conformal factor (and its spatial derivatives)
(use\_conformal=''yes''). In general, the option use\_conformal=''yes'' should 
be used, since ????. 

\section{Parameters}

\subsection{List of all Parameters}

\subsection{Discussion}

\section{Interaction with Other Thorns}

It is still to be decided how initial data should be supplied to
for systems of evolution equations which use more or different 
variables to those in the einstein implementation

\section{Future Development}

Initial data sets which are missing from this thorn include

\begin{itemize}

\item{} Boosted single black hole
\item{} Single black hole in harmonic spatial coordinates
\item{} Single black hole in Eddington-Finkelstein coordinates

\end{itemize}


\end{document}
